We compare denoising diffusion (DDPM with DDIM and ancestral samplers), conditional flow matching with Euler and Heun samplers, and one-round rectified-flow reflow on three 2D toy densities, using one shared MLP, one training budget, five seeds, and sliced Wasserstein and MMD scores against 10000 held-out points for 1 to 100 sampling steps. All methods are reimplemented. Flow matching has lower one-step sliced Wasserstein (SW) than DDIM on every task but higher one-step MMD, and neither is usable at one step. At four steps it has lower SW on eight Gaussians, marginally lower SW on the checkerboard, and \emph{higher} SW on two moons; this reversal is SW-only, since MMD ranks flow matching ahead of DDIM at four steps on all three tasks. Whether its advantage shrinks with more steps is mixed: the SW ratio to DDIM shrinks on eight Gaussians, grows on the checkerboard and changes side on two moons. Flow-matching trajectories are measured as \emph{less} straight than DDIM trajectories by chord-to-arc length. One reflow round brings one-step SW close to the real-data floor (\rhval{main/reflow-1-euler/eight_gaussians/sw_k1/mean:3} against \rhval{main/real-data-floor/eight_gaussians/sw_k1/mean:3} on eight Gaussians, where the gap is widest) and gives no benefit at 100 steps; the student inherits the quality of its teacher's sampler (a 4-step teacher limits it at every reported step count), and a second round adds little. Without retraining, a Heun sampler on the flow-matching model is better than Euler at 8 network evaluations on two moons and the checkerboard and worse at 2. In an ablation, cosine-schedule DDIM fails at every tested step count unless its sampling grid is started below the last timestep. The study is small (2D, one network size, 3000 updates, no tuning), and the reflow comparison is not matched in training compute.
